\chapter{Quantum Continuity Equation}

\section*{Introduction}

This equation maps the abstract time evolution of the wavefunction onto the concrete physical requirement that probability is neither created nor destroyed. If the probability density increases somewhere, it must be because probability has flowed in from elsewhere---and the continuity equation precisely quantifies this balance.


\begin{equation}
\frac{\partial |\psi|^2}{\partial t} + \nabla \cdot \mathbf{J} = 0
\end{equation}

\section*{Fundamental Equations Used}

The derivation starts from the time-dependent Schr\"odinger equation and its complex conjugate.

\section*{Assumptions}

\textbf{Assumptions:} Non-relativistic quantum mechanics; the Hamiltonian is Hermitian; the wavefunction is normalizable.


\textbf{Limitations:} Does not apply when particle creation or annihilation occurs (quantum field theory); requires modification for open quantum systems with decoherence.


\section*{Mathematical Derivation}

\textbf{Step 1: Write the time-dependent Schr\"{o}dinger equation and its conjugate.}

The Schr\"{o}dinger equation for a single particle in a potential $V$ is
\begin{equation}
i\hbar\frac{\partial\psi}{\partial t} = -\frac{\hbar^2}{2m}\nabla^2\psi + V\psi.
\end{equation}
Taking the complex conjugate (noting $V$ is real):
\begin{equation}
-i\hbar\frac{\partial\psi^*}{\partial t} = -\frac{\hbar^2}{2m}\nabla^2\psi^* + V\psi^*.
\end{equation}

\textbf{Step 2: Multiply by $\psi^*$ and $\psi$ respectively.}

\begin{align}
i\hbar\,\psi^*\frac{\partial\psi}{\partial t} &= -\frac{\hbar^2}{2m}\,\psi^*\nabla^2\psi + V|\psi|^2, \\
-i\hbar\,\psi\frac{\partial\psi^*}{\partial t} &= -\frac{\hbar^2}{2m}\,\psi\,\nabla^2\psi^* + V|\psi|^2.
\end{align}

\textbf{Step 3: Subtract the second equation from the first.}

The potential terms cancel:
\begin{equation}
i\hbar\left(\psi^*\frac{\partial\psi}{\partial t} + \psi\frac{\partial\psi^*}{\partial t}\right) = -\frac{\hbar^2}{2m}\left(\psi^*\nabla^2\psi - \psi\,\nabla^2\psi^*\right).
\end{equation}
The left-hand side is $i\hbar\,\partial|\psi|^2/\partial t$.

\textbf{Step 4: Express the right-hand side as a divergence.}

Using $\psi^*\nabla^2\psi - \psi\,\nabla^2\psi^* = \nabla\cdot(\psi^*\nabla\psi - \psi\,\nabla\psi^*)$:
\begin{equation}
i\hbar\frac{\partial|\psi|^2}{\partial t} = -\frac{\hbar^2}{2m}\nabla\cdot\left(\psi^*\nabla\psi - \psi\,\nabla\psi^*\right).
\end{equation}

\textbf{Step 5: Divide by $i\hbar$ and identify the probability current.}

\begin{equation}
\frac{\partial|\psi|^2}{\partial t} = -\frac{\hbar}{2mi}\nabla\cdot\left(\psi^*\nabla\psi - \psi\,\nabla\psi^*\right) = -\nabla\cdot\mathbf{J},
\end{equation}
where $\mathbf{J} = \frac{\hbar}{2mi}(\psi^*\nabla\psi - \psi\,\nabla\psi^*)$ is the probability current. Thus we arrive at
\begin{equation}
\boxed{\frac{\partial|\psi|^2}{\partial t} + \nabla\cdot\mathbf{J} = 0.}
\end{equation}
This is the quantum continuity equation, expressing local conservation of probability.

\section*{Characteristics of the Equation}

\begin{itemize}
    \item This is a local conservation law---probability is conserved point by point, not just globally.
    \item Derived directly from the Schrodinger equation and its complex conjugate.
    \item The integrated form $\frac{d}{dt}\int|\psi|^2 dV = 0$ follows by the divergence theorem.
    \item It is the quantum analog of the continuity equation in fluid dynamics and electrodynamics.
\end{itemize}
\section*{Solution Methods}

\begin{itemize}
    \item \textbf{Derivation from Schrodinger equation:} Multiply the Schrodinger equation by $\psi^*$, the conjugate by $\psi$, and subtract to obtain the continuity equation.
    \item \textbf{Verification by substitution:} Given a specific $\psi(x,t)$, verify that the equation holds.
    \item \textbf{Conservation check:} Compute $\frac{d}{dt}\int|\psi|^2 dV$ and show it vanishes using integration by parts.
\end{itemize}
\section*{Verifications}

\begin{itemize}
    \item For any normalizable wavefunction evolving under a Hermitian Hamiltonian, $\int|\psi|^2 dV = 1$ for all $t$.
    \item For a free-particle Gaussian wavepacket, the spreading is precisely compensated by the inward/outward probability flux.
    \item In scattering problems, the incoming and outgoing probability fluxes are equal, ensuring unitarity.
\end{itemize}
\section*{Applications}

\begin{itemize}
    \item Proof that the total probability remains unity for all time.
    \item Establishing the probability current as a physically meaningful quantity.
    \item Foundation for charge conservation in quantum electrodynamics.
    \item Ensuring consistency of time-dependent numerical simulations.
\end{itemize}
\section*{What Richard Feynman Would Have Asked}

\subsection*{Plain-English Translation}
\textit{``Can you explain the quantum continuity equation in a way that a first-year student would understand, without using any equations?''}

Probability is like water: it cannot appear or disappear from anywhere. If the amount of ``probability water'' in a bucket increases, it must have flowed in from somewhere else through the walls of the bucket. The quantum continuity equation is exactly this accounting statement, applied to every tiny region of space simultaneously. It is the reason quantum mechanics is self-consistent: the total probability of finding the particle anywhere is always exactly one.

\subsection*{Localized Mechanics}
\textit{``What exactly is happening at a point where the probability density is increasing? Where does the new probability come from?''}

When the probability density increases at a point, it means the divergence of the probability current is negative at that point---probability is flowing inward from all directions, converging on that location. Conversely, where the probability density decreases, the current diverges outward, carrying probability away. The continuity equation is the local statement that these two effects balance perfectly: the rate of accumulation equals the rate of convergence. No probability is created or destroyed; it merely shifts from place to place.

\subsection*{Testing Extreme Limits}
\textit{``What happens in the extreme limit where a particle is completely localized---a delta function? Does the continuity equation still hold?''}

A delta function is not a valid wavefunction in the traditional sense (it is not square-integrable), but as a limiting case, the continuity equation still applies in the distributional sense. A perfect localization requires infinite momentum uncertainty, and the immediate spreading of the probability (a delta function ``explodes'' into a broad distribution) is precisely what the continuity equation predicts: the outward probability current from the point of localization is enormous, causing rapid depletion at the center and growth at the edges. The equation handles this extreme limit correctly.

\subsection*{Dimensionality and Geometry}
\textit{``Does the continuity equation look different in one, two, and three dimensions, and does the dimensionality affect the physics?''}

The mathematical form $\partial\rho/\partial t + \nabla \cdot \mathbf{J} = 0$ is identical in all dimensions---only the divergence operator changes form. In one dimension it becomes $\partial\rho/\partial t + \partial J_x/\partial x = 0$, which is even simpler. The dimensionality does not change the physics of probability conservation, but it does affect the types of solutions: in one dimension, probability can only flow left or right, while in three dimensions it can flow in any direction, leading to much richer dynamics such as scattering in multiple directions.

\subsection*{Cross-Disciplinary Connection}
\textit{``This continuity equation looks exactly like one I saw in fluid dynamics or electromagnetism. Is it literally the same equation?''}

Yes, it is structurally identical to the continuity equation for mass in fluid dynamics ($\partial\rho/\partial t + \nabla \cdot (\rho\mathbf{v}) = 0$) and for charge in electromagnetism ($\partial\rho/\partial t + \nabla \cdot \mathbf{J} = 0$). The only difference is the physical interpretation of the density and current. This universality is not a coincidence---it reflects the deep mathematical structure of conservation laws. Any conserved quantity in physics must satisfy a continuity equation of this form, and the quantum continuity equation is simply the instance where the conserved quantity is probability.

% ------------------------------------------------------------
\section*{FAQ}

\textbf{Q: Does the continuity equation hold for time-independent problems?}
A: Yes, but trivially. For stationary states, $\partial|\psi|^2/\partial t = 0$, so $\nabla \cdot \mathbf{J} = 0$---the probability current is divergenceless, meaning there is no net accumulation or depletion of probability anywhere.

\textbf{Q: What if the Hamiltonian is not Hermitian?}
A: If the Hamiltonian is non-Hermitian (as in effective descriptions of open systems with gain or loss), the continuity equation acquires a source term. Probability is no longer conserved, reflecting the coupling to an external environment.

\textbf{Q: Is this related to charge conservation in electromagnetism?}
A: Yes, directly. The electromagnetic continuity equation $\partial\rho/\partial t + \nabla \cdot \mathbf{J} = 0$ is the classical precursor. The quantum continuity equation is its non-relativistic quantum counterpart applied to probability rather than charge.
\section*{Important Bibliography}

\begin{enumerate}
\item Griffiths, D.J., \textit{Introduction to Quantum Mechanics}, 3rd ed. (Cambridge University Press, 2018), Chapter 1.
\item Shankar, R., \textit{Principles of Quantum Mechanics}, 2nd ed. (Springer, 1994), Chapter 5.
\item Sakurai, J.J. and Napolitano, J., \textit{Modern Quantum Mechanics}, 2nd ed. (Cambridge University Press, 2017), Chapter 1.
\end{enumerate}

